\documentclass[10pt,letterpaper]{article}
\usepackage[margin=1in]{geometry}
\usepackage{fontspec}
\usepackage{microtype}
\usepackage{xcolor}
\usepackage[hidelinks]{hyperref}
\usepackage{parskip}
\setmainfont{TeX Gyre Termes}
\setsansfont{TeX Gyre Heros}
\setlength{\parindent}{0pt}
\setlength{\parskip}{6.2pt}
\pagestyle{empty}
\hypersetup{
  pdftitle={Cover letter to Nature Communications},
  pdfauthor={Huan Zhong}
}

\begin{document}
{\sffamily\bfseries\Large Nature Communications}\par
\vspace{2pt}\hrule height 0.7pt\vspace{15pt}
\begin{flushright}12 November 2026\end{flushright}
Editors\\
\textit{Nature Communications}

\textbf{Re: WrapEvoFS enables auditable feature compression with regret-constrained representative locking}

Dear Editors,

We submit the Article entitled ``WrapEvoFS enables auditable feature compression with regret-constrained representative locking'' for consideration in \textit{Nature Communications}.

Repeated stochastic searches can produce different near-scoring feature signatures, yet one must often be retained for a model or follow-up analysis. WrapEvoFS makes that decision explicit: it restricts candidates to an explicit empirical development-score tolerance, selects the feature-overlap medoid within that pool, and records the chosen signature before held-out evaluation.

In a fully nested four-class TCGA analysis of 491 participants, the workflow reduced screened feature counts by 80--82\%, retaining means of 15--20 features with macro AUROC estimates of 0.827--0.828. Absolute AUROC costs relative to Direct screening were 0.004--0.016. Under the fixed reference protocol, the rule changed 11 of 30 highest-score selections and increased eligible-pool overlap. Top-three and unrestricted medoids exceeded that score gap in 5 and 14 banks, respectively. Native and size-matched selector controls quantify the performance--compression trade-offs across branches.

A retrospective evaluation of all 150 saved candidates further illustrates why the reporting rule should precede outcome inspection. Across repeated half-splits of the outer-test data, selecting the best run on one half gave an average selection-minus-opposite-half AUROC difference of 0.0253. This conditional diagnostic complements a one-time evaluation of 24 frozen signatures across four held-out AMP-AD centers and controlled simulations identifying when informative overlap can help resolve noisy score rankings.

WrapEvoFS is available under the BSD 3-Clause License with tagged releases and PyPI distributions. Reference analyses retain v0.2.0, while v0.3.0 adds the opt-in calibration API. Source, analysis code, aggregate source data, checksums, and provenance records are available from the \href{https://github.com/junseoparkX/wrapevofs-manuscript-reproducibility}{manuscript reproducibility repository}. Participant-level access routes and restrictions are stated in the manuscript.

The manuscript is original, is not under consideration elsewhere, and has not been discussed previously with a \textit{Nature Communications} editor. Authorship, contributions, funding, and competing-interest declarations are provided in the manuscript. Reviewer suggestions and exclusions will be entered in the submission system.

The submission includes five main figures and two main tables, with 29 Supplementary Figures and 33 Supplementary Tables.

Thank you for considering our work.

Sincerely,\\[3pt]
\textbf{Huan Zhong}\\
Michael Smith Laboratories\\
Department of Biochemistry and Molecular Biology\\
University of British Columbia\\
Vancouver, BC, Canada\\
\href{mailto:huan.zhong@ubc.ca}{huan.zhong@ubc.ca}
\end{document}
